\documentclass[11pt]{article}
\usepackage[utf8]{inputenc}
\usepackage{amsmath,amssymb,amsthm}
\usepackage[margin=1in]{geometry}
\usepackage{hyperref}
\usepackage{xcolor}
\usepackage{enumitem}
\usepackage{newunicodechar}
\newunicodechar{ℝ}{\ensuremath{\mathbb{R}}}
\newunicodechar{ℕ}{\ensuremath{\mathbb{N}}}
\newunicodechar{α}{\ensuremath{\alpha}}

\title{Tide retrospective:\\\emph{Mixed odd-power 2D moment vanishing (V1)}}
\author{Daniel Murfet}
\date{18 May 2026}

\begin{document}
\maketitle

\begin{abstract}
This retrospective records a small parity-completion pair on top of
the quartic--sextic 2D potential machinery established by Tide~A2
and used by Tide~T1. For $L(x, y) = x^4/24 + y^6/720$ at $t > 0$,
the two theorems state that any joint moment
$\langle x^a y^b \rangle_{2D, t}$ vanishes when at least one of
$a, b$ is odd, with no case-split: by leaving the non-odd index as a
free natural number rather than forcing it to $2k$, the pair of
theorems together covers the odd-odd case as well. Each proof is a
single \texttt{rw} block (factor theorem + 1D odd-vanishing) closed
by \texttt{simp}; the file is 80 lines including docstrings. Single
session, build green on first compile after the integrability and
1D lemmas were located. Zero \texttt{sorry}, zero \texttt{axiom},
zero \texttt{native\_decide}.
\end{abstract}

\section{Setting}

The grammar-paper precursor sequence and the 2D quartic--sextic
sub-thread of Tide~A2 left T1 as the first joint-moment theorem on
the new 2D potential: the even-power closed form
$\langle x^{2j} y^{2k} \rangle_{2D, t}$. T1's retrospective
explicitly flagged the parity-completion pair as a deferred
follow-up (T1's ``Candidate~C''): the mirror odd-vanishing theorems
needed to give a full picture of mixed-monomial moments under the
separable measure. The 1D building blocks\footnote{Specifically
\texttt{Laplace.OneD.quartic\_expected\_value\_odd} and
\texttt{Laplace.OneD.sextic\_expected\_value\_odd}.} were already
in the seabed, both proven unconditionally in $t$ via odd-symmetry
of the integrand against the even-symmetric exponential factor.
A2's factor theorem\footnote{\texttt{Laplace.TwoD.gibbsExpectation\_quarticSextic\_factor},
in \texttt{Laplace/TwoD/QuarticSextic.lean}.} provides the
2D-to-1D reduction. With all three pieces in place, V1 is the
minimal step that finishes the parity story.

\section{The thing formalised}

For the 2D potential $L(x, y) = x^4/24 + y^6/720$ at temperature
$t > 0$:

\begin{itemize}[leftmargin=1.5em]
\item For all $j, n \in \mathbb N$,
  $\langle x^{2j+1} y^n \rangle_{2D, t} = 0$.
\item For all $m, k \in \mathbb N$,
  $\langle x^m y^{2k+1} \rangle_{2D, t} = 0$.
\end{itemize}

The two theorems together cover every $\langle x^a y^b \rangle_{2D, t}$
with at least one of $a, b$ odd, including the odd-odd case. Lean
signatures:

\begin{verbatim}
theorem gibbsExpectation_quarticSextic_odd_pow_pow_eq_zero
    (j n : ℕ) {t : ℝ} (ht : 0 < t) :
    gibbsExpectation
        (addSeparable Laplace.OneD.quarticPotential
                      Laplace.OneD.sexticPotential) t
        (fun z : ℝ × ℝ => z.1 ^ (2 * j + 1) * z.2 ^ n) = 0

theorem gibbsExpectation_quarticSextic_pow_odd_pow_eq_zero
    (m k : ℕ) {t : ℝ} (ht : 0 < t) :
    gibbsExpectation
        (addSeparable Laplace.OneD.quarticPotential
                      Laplace.OneD.sexticPotential) t
        (fun z : ℝ × ℝ => z.1 ^ m * z.2 ^ (2 * k + 1)) = 0
\end{verbatim}

Both live in the \texttt{Laplace.TwoD} namespace, conforming to
T1's even-case naming.\footnote{T1's headline theorem is
\texttt{gibbsExpectation\_quarticSextic\_pow\_pow\_eq}; the new pair
mirrors that name with an \texttt{odd}-infix and \texttt{\_eq\_zero}
suffix.} The new file is
\texttt{Laplace/TwoD/QuarticSexticOddMoments.lean}, sibling to T1's
\texttt{QuarticSexticMoments.lean}, and is added to the project root
\texttt{Laplace.lean}.

\section{Proof strategy}

For the odd-$x$ side: apply the factor theorem with
$f = x \mapsto x^{2j+1}$ and $g = y \mapsto y^n$. The 2D expectation
factors as a product of the 1D quartic expectation and the 1D sextic
expectation. The first factor is zero by
\texttt{quartic\_expected\_value\_odd}. The product is zero, and
\texttt{simp} closes the residual \texttt{0 * \_} via
\texttt{zero\_mul}. The odd-$y$ side is the mirror image: the
\emph{second} factor is zero by \texttt{sextic\_expected\_value\_odd},
\texttt{simp} closes via \texttt{mul\_zero}.

The integrability hypotheses fed to the factor theorem are the
canonical \texttt{\_integrable\_pow\_pot} witnesses\footnote{Namely
\texttt{Laplace.OneD.quartic\_integrable\_pow\_pot} and
\texttt{Laplace.OneD.sextic\_integrable\_pow\_pot}.}, the same shape
T1 used. They unify with the factor theorem's expected form
definitionally; no \texttt{change}-tactic bridging needed.

\section{Roadblocks and resolutions}

There were no roadblocks. The proof template was correct on first
write; \texttt{lake build} returned green on the first compile. The
deliberation surfaced one strict-improvement over the user's draft
(see \S\ref{sec:gpt}), so the implementation step ran on the
deliberation's revised plan rather than the user's initial sketch.

\section{What was Mathlib, what was new}\label{sec:gpt}

\paragraph{From Mathlib.} Nothing direct in this tide. The factor
theorem and the 1D odd-vanishing theorems are Laplace-side; the
only Mathlib-side dependency is \texttt{zero\_mul} / \texttt{mul\_zero}
inside \texttt{simp}.

\paragraph{From the seabed.} Three pieces did the work: A2's
factorisation engine, the 1D quartic and sextic odd-vanishing
theorems (both unconditional in $t$), and the canonical
integrability lemmas. Names recorded in the footnotes of the
\S\emph{Setting} and \S\emph{Proof strategy} sections above.

\paragraph{New here.} Two parity-completion theorems, ~80 lines total
(including docstrings). The architectural contribution is the choice
to leave the non-odd index as a free $n$ (resp.\ $m$) rather than
forcing it to $2k$ (resp.\ $2j$): this is what makes the pair cover
the odd-odd case without a case split or an \texttt{Odd}-witness
predicate. The contribution came from GPT-5.5~Pro at Step~2; see the
Lessons section.

\section{Lessons}

\paragraph{Free-the-other-index is a strict improvement on
forced-even.} The user's initial draft had the two theorems keyed on
$(2j+1, 2k)$ and $(2j, 2k+1)$, mirroring T1's even-case
$(2j, 2k)$ structure. GPT-5.5~Pro pointed out this misses the
odd-odd case and recommended $(2j+1, n)$ and $(m, 2k+1)$: same proof
size, no case split, total coverage. The pattern generalises:
\emph{when proving a vanishing identity that factors through a single
factor's vanishing, leave the other factor's parameter free rather
than mirror the parametrisation of the matching positive identity.}
The matching positive identity (T1) needs both factors to have
closed forms, hence the symmetric $(2j, 2k)$; the vanishing identity
(V1) only needs one factor to vanish, so the symmetry breaks and the
freer parametrisation costs nothing.

\paragraph{\texttt{simp} over \texttt{ring} for residual
\texttt{0 \(*\) \_} goals.} GPT also suggested \texttt{simp} in
place of the draft's \texttt{ring} at the bottom of each proof.
After the two rewrites, the goal is literally \texttt{0 * \_ = 0}
or \texttt{\_ * 0 = 0}; \texttt{simp} handles this via
\texttt{zero\_mul} / \texttt{mul\_zero} (and \texttt{simp} is at this
point cheaper than \texttt{ring}, which would also work but would
go through a normalisation pass). Worth noting as a tactic
preference for this class of small wrap-up goal.

\paragraph{The single-round consensus.} Both Claude and GPT-5.5~Pro
agreed on Candidate~A$'$ in one round. The vote, the integrability
unification, the bottom-tactic choice, and the rejection of the
generic \texttt{addSeparable} abstraction (``one tide too far'')
all landed in the same consult. Pattern: small tides off a recent
retrospective's flagged Candidate often have one canonically-right
shape that a single GPT round nails. The cost of the consult was
recouped by the strict improvement on the user's draft.

\section{Follow-ups}

\begin{itemize}[leftmargin=1.5em]
\item \emph{Generic \texttt{addSeparable} odd-vanishing abstraction.}
  For an arbitrary additively separable 2D potential
  $V(x, y) = V_1(x) + V_2(y)$ where $V_1$ is even-symmetric, any
  monomial moment with an odd $x$-exponent vanishes. The 1D abstract
  prerequisite is ``$\int x^{2j+1} e^{-tV(x)}\,dx = 0$ for $V$
  even-symmetric'' --- which the seabed does not yet have as a
  generic lemma (only the quartic and sextic specialisations). When a
  third or fourth even-symmetric 1D potential lands (e.g.\ a pure
  $k$-th-power potential beyond $k = 4, 6$), promote the generic
  abstraction and rewrite the quartic and sextic odd-vanishings as
  specialisations.
\item \emph{Higher-power named consumers.} If a downstream consumer
  ever wants a specific named family
  $\langle x^{2j+1} y^{2k}\rangle_{2D, t} = 0$ (rather than the
  generic $(2j+1, n)$ form), it is a one-line corollary of V1's
  first theorem at the call site; no follow-up tide needed.
\item \emph{\texttt{\textbackslash leanref}-tag the primer's 2D
  examples} once the primer's TeX has statements for the 2D
  parity-completion identities. Pin to this tide's commit. Pairs with
  the existing K2 / U3 / L5 paper-tagging backlog.
\end{itemize}

\section*{Acknowledgements}

GPT-5.5~Pro's deliberation contributed three load-bearing pieces of
guidance in a single round: the parametrisation strengthening from
$(2j+1, 2k)$ to $(2j+1, n)$ that closes the odd-odd hole the draft
missed; \texttt{simp} as the bottom tactic in place of \texttt{ring};
and the \texttt{\_eq\_zero} naming suffix. The proof landed
single-attempt with no thrashing-rescue consult; the full GPT
response is preserved alongside the tide log in the primer project's
tide-log folder.

This Tide branched off the post-L3-merge \texttt{main}, the same
\texttt{main} that gathered the seven-tide backlog merged earlier on
2026-05-18 (the day's bookkeeping cleanup). The fresh \texttt{main}
contains A2, T1, U2, R1, Q1, L2, L3 inline rather than as detached
tide branches.

\end{document}
